\documentclass[11pt,a4paper]{article}

% Packages
\usepackage{amsmath,amssymb,amsthm}
\usepackage{mathtools}
\usepackage{physics}
\usepackage{graphicx}
\usepackage{geometry}
\usepackage{fancyhdr}
\usepackage{hyperref}
\usepackage{xcolor}
\usepackage{tcolorbox}
\usepackage{enumitem}
\usepackage{booktabs}

% Page geometry
\geometry{margin=1in}

% Header/Footer
\pagestyle{fancy}
\fancyhf{}
\fancyhead[L]{\small Derivation of the Chalmers Ansatz}
\fancyhead[R]{\small K-Dense Web}
\fancyfoot[C]{\thepage}
\fancyfoot[R]{\small Generated using K-Dense Web (\href{https://k-dense.ai}{k-dense.ai})}

% Theorem environments
\theoremstyle{definition}
\newtheorem{definition}{Definition}[section]
\newtheorem{axiom}{Axiom}[section]
\newtheorem{principle}{Principle}[section]
\theoremstyle{plain}
\newtheorem{theorem}{Theorem}[section]
\newtheorem{lemma}{Lemma}[section]
\newtheorem{corollary}{Corollary}[section]
\theoremstyle{remark}
\newtheorem{remark}{Remark}[section]

% Custom commands
\newcommand{\logent}{\vDash_{\text{log}}}
\newcommand{\notlogent}{\nvDash_{\text{log}}}
\newcommand{\natent}{\vDash_{\text{nat}}}
\newcommand{\Wlog}{W_{\text{log}}}
\newcommand{\Wnat}{W_{\text{nat}}}
\newcommand{\Forg}{F_{\text{org}}}
\newcommand{\PhiP}{\Phi_{P}}
\newcommand{\PhiE}{\Phi_{E}}
\newcommand{\Pactual}{P_{\text{actual}}}
\newcommand{\Psiuniv}{\Psi_{\text{univ}}}

\title{\textbf{Derivation of the Chalmers Ansatz}\\[0.5em]
\large A Formal Mathematical Framework for Naturalistic Dualism}
\author{K-Dense Web\\
\small \href{mailto:contact@k-dense.ai}{contact@k-dense.ai}}
\date{February 2026}

\begin{document}

\maketitle

\begin{abstract}
This document provides a rigorous symbolic formalization of David Chalmers' Naturalistic Dualism as presented in \textit{The Conscious Mind}. We define a tuple $\mathcal{U} = \langle P, M, E, L \rangle$ representing the universe, where $P$ is the causally closed physical domain, $M$ is the psychological/functional domain, $E$ is the phenomenal domain (qualia), and $L$ is the set of fundamental laws. Using standard Kripke semantics, we formalize the Zombie Argument to prove that logical supervenience fails for consciousness ($P \notlogent E$), necessitating fundamental psychophysical laws. We derive the mapping function $\Psi: P \to E$ for natural supervenience, the Principle of Structural Coherence ($S(E) \cong S(A)$), the Principle of Organizational Invariance, and the Double-Aspect Principle of Information. Finally, we extend the framework to quantum mechanics, showing how superposed physical states yield distinct phenomenal subjects.
\end{abstract}

\begin{figure}[h!]
\centering
\includegraphics[width=0.85\textwidth]{../figures/graphical_abstract.png}
\caption{\textbf{Graphical Abstract.} The dual-aspect framework of Naturalistic Dualism. The Physical Domain $P$ (causally closed) relates to the Phenomenal Domain $E$ via Natural Supervenience ($\Psi$), not Logical Supervenience. The Psychological Domain $M$ logically supervenes on $P$. Information Space $I$ serves as the ontological bridge with dual realizations $\PhiP$ and $\PhiE$.}
\label{fig:graphical_abstract}
\end{figure}

\newpage

%==============================================================================
\section{Legend of Variables}
%==============================================================================

\begin{tcolorbox}[colback=gray!5,colframe=gray!50,title=\textbf{Primary Domains}]
\begin{tabular}{@{}ll@{}}
$P$ & Physical domain: the set of all physical states/facts \\
$M$ & Psychological domain: functional/causal-role states (judgments, dispositions) \\
$E$ & Phenomenal domain: the set of all experiential/qualitative states (qualia) \\
$L$ & Set of fundamental laws: $L = L_P \cup L_\Psi$ \\
$\mathcal{U}$ & Universe tuple: $\mathcal{U} = \langle P, M, E, L \rangle$
\end{tabular}
\end{tcolorbox}

\begin{tcolorbox}[colback=blue!5,colframe=blue!50,title=\textbf{Modal \& Semantic Operators}]
\begin{tabular}{@{}ll@{}}
$\Wlog$ & Set of all logically possible worlds \\
$\Wnat$ & Set of all naturally (nomologically) possible worlds; $\Wnat \subset \Wlog$ \\
$w$ & A possible world; $w \in \Wlog$ \\
$P(w)$ & Physical instantiation in world $w$ \\
$E(w)$ & Phenomenal instantiation in world $w$; $E(w) = \varnothing$ denotes absence \\
$\logent$ & Logical supervenience/entailment: $A \logent B$ iff $\forall w \in \Wlog: A(w) \Rightarrow B(w)$ \\
$\natent$ & Natural supervenience: $A \natent B$ iff $\forall w \in \Wnat: A(w) \Rightarrow B(w)$
\end{tabular}
\end{tcolorbox}

\begin{tcolorbox}[colback=green!5,colframe=green!50,title=\textbf{Mapping Functions}]
\begin{tabular}{@{}ll@{}}
$\Psi$ & Psychophysical bridging function: $\Psi: P \to E$ (natural supervenience) \\
$\Forg(x)$ & Fine-grained functional organization of system $x$ \\
$S(\cdot)$ & Structure operator extracting relational/geometric structure \\
$A$ & Awareness: the subset of $M$ directly available for global control \\
$I$ & Information space: the set of all informational states
\end{tabular}
\end{tcolorbox}

\begin{tcolorbox}[colback=orange!5,colframe=orange!50,title=\textbf{Information \& Realization}]
\begin{tabular}{@{}ll@{}}
$\PhiP$ & Physical realization function: $\PhiP: I \to P$ \\
$\PhiE$ & Phenomenal realization function: $\PhiE: I \to E$ \\
$s \in I$ & An informational state
\end{tabular}
\end{tcolorbox}

\begin{tcolorbox}[colback=purple!5,colframe=purple!50,title=\textbf{Quantum Mechanical Notation}]
\begin{tabular}{@{}ll@{}}
$\ket{\Psiuniv}$ & Universal quantum state \\
$\ket{P_i}$ & Basis state of physical implementation $i$ \\
$c_i \in \mathbb{C}$ & Complex amplitude; $\sum_i |c_i|^2 = 1$ \\
$E_i$ & Phenomenal subject associated with branch $i$
\end{tabular}
\end{tcolorbox}

\newpage

%==============================================================================
\section{Derivation of the Chalmers Ansatz}
%==============================================================================

%------------------------------------------------------------------------------
\subsection{Fundamental Domains (Chapters 1 \& 4)}
%------------------------------------------------------------------------------

\begin{definition}[Universe Tuple]
The Chalmers universe is a 4-tuple:
\begin{equation}
\boxed{\mathcal{U} = \langle P, M, E, L \rangle}
\end{equation}
where $P$, $M$, $E$ are pairwise distinct domains and $L = L_P \cup L_\Psi$ comprises physical and psychophysical laws.
\end{definition}

\begin{axiom}[Causal Closure of the Physical]\label{ax:causal_closure}
The physical domain $P$ is causally closed:
\begin{equation}
\forall p, p' \in P: \text{Cause}(p, p') \Rightarrow p' \in P
\end{equation}
All physical effects have sufficient physical causes. The dynamics of $P$ are governed exclusively by $L_P$.
\end{axiom}

\begin{definition}[Psychological Domain]
The psychological domain $M$ consists of functional states characterized by causal role:
\begin{equation}
M \coloneqq \{ m : m \text{ is individuated by input-output relations} \}
\end{equation}
Elements of $M$ include beliefs, desires, learning states, and judgments---defined dispositionally.
\end{definition}

\begin{theorem}[Logical Supervenience of $M$ on $P$]\label{thm:M_supervenes}
The psychological domain logically supervenes on the physical:
\begin{equation}
\boxed{P \logent M}
\end{equation}
\end{theorem}

\begin{proof}
Let $w_1, w_2 \in \Wlog$ with $P(w_1) = P(w_2)$. Since $M$ is functionally defined (causal-role states), and functional organization is determined by physical structure, we have $M(w_1) = M(w_2)$. Hence $P \logent M$. \qed
\end{proof}

\begin{definition}[Phenomenal Domain]
The phenomenal domain $E$ is the set of all qualitative experiential states:
\begin{equation}
E \coloneqq \{ e : e \text{ is characterized by ``what it is like'' to instantiate } e \}
\end{equation}
Elements of $E$ are qualia: the redness of red, the painfulness of pain, etc.
\end{definition}

\begin{axiom}[Fundamentality of $E$]\label{ax:fundamental_E}
The phenomenal domain $E$ is a fundamental feature of $\mathcal{U}$, irreducible to $P$ or $M$:
\begin{equation}
E \not\subseteq P \quad \text{and} \quad E \not\subseteq M
\end{equation}
Phenomenal properties are ontologically primitive.
\end{axiom}

%------------------------------------------------------------------------------
\subsection{The Hard Problem and Failure of Reduction (Chapter 3)}
%------------------------------------------------------------------------------

We now formalize the Zombie Argument using Kripke semantics to prove that logical supervenience fails for consciousness.

\begin{definition}[Kripke Frame for Modal Analysis]
Let $\mathcal{K} = (\Wlog, R)$ be a Kripke frame where:
\begin{itemize}[nosep]
    \item $\Wlog$ is the set of all logically possible worlds
    \item $R \subseteq \Wlog \times \Wlog$ is the accessibility relation (for logical possibility, $R = \Wlog \times \Wlog$)
\end{itemize}
For logical possibility, $\Diamond \phi$ holds at $w$ iff $\exists w' \in \Wlog: \phi$ holds at $w'$.
\end{definition}

\begin{definition}[Zombie World]
A zombie world $w_z \in \Wlog$ is defined by:
\begin{equation}
P(w_z) = \Pactual \quad \land \quad E(w_z) = \varnothing
\end{equation}
where $\Pactual$ denotes the complete physical state of the actual world.
\end{definition}

\begin{theorem}[Zombie Argument---Failure of Logical Supervenience]\label{thm:zombie}
There exists a logically possible world physically identical to the actual world but lacking phenomenal experience:
\begin{equation}
\boxed{\exists w \in \Wlog : \big( P(w) = \Pactual \big) \land \big( E(w) = \varnothing \big)}
\end{equation}
\end{theorem}

\begin{proof}
\textbf{Step 1 (Conceivability):} The concept of a zombie (a being physically identical to a conscious being but lacking experience) involves no logical contradiction. The primary intension of ``physical'' and ``phenomenal'' are conceptually independent.

\textbf{Step 2 (Conceivability-Possibility Link):} For ideal conceivability (not ruled out a priori upon sufficient reflection), if $\phi$ is ideally conceivable, then $\Diamond \phi$ (logical possibility).

\textbf{Step 3 (Zombie Conceivability):} A zombie world is ideally conceivable: complete knowledge of physics does not entail knowledge of experience. No amount of physical information logically necessitates phenomenal facts.

\textbf{Step 4 (Conclusion):} Therefore, $\exists w_z \in \Wlog$ satisfying the zombie conditions. \qed
\end{proof}

\begin{corollary}[Failure of Logical Supervenience for $E$]\label{cor:no_log_supervenience}
Phenomenal properties do not logically supervene on physical properties:
\begin{equation}
\boxed{P \notlogent E}
\end{equation}
\end{corollary}

\begin{proof}
Direct from Theorem~\ref{thm:zombie}. If $P \logent E$ held, then $\forall w \in \Wlog: P(w) = \Pactual \Rightarrow E(w) = E_{\text{actual}}$. But $w_z$ is a counterexample. \qed
\end{proof}

\begin{corollary}[Materialism is False]
Since $P \notlogent E$ and consciousness exists, materialism (the thesis that all facts globally logically supervene on physical facts) is false:
\begin{equation}
\neg \big( \forall F: P \logent F \big)
\end{equation}
\end{corollary}

%------------------------------------------------------------------------------
\subsection{Psychophysical Laws (Chapters 4 \& 6)}
%------------------------------------------------------------------------------

Given that $P \notlogent E$, new fundamental laws are required.

\begin{axiom}[Natural Supervenience of $E$ on $P$]\label{ax:natural_supervenience}
Phenomenal properties naturally supervene on physical properties:
\begin{equation}
P \natent E
\end{equation}
For all nomologically possible worlds: $\forall w \in \Wnat: P(w) = P(w') \Rightarrow E(w) = E(w')$.
\end{axiom}

\begin{definition}[Psychophysical Bridging Function $\Psi$]
The fundamental psychophysical law is encoded in a mapping:
\begin{equation}
\boxed{\Psi: P \to E}
\end{equation}
such that for any physical state $p \in P$, $\Psi(p) \in E$ is the associated phenomenal state. The function $\Psi$ represents nomological (natural) necessity, distinct from logical necessity.
\end{definition}

\begin{remark}
$\Psi$ is not derivable from $L_P$; it is a fundamental addition to the laws of nature, analogous to the introduction of electromagnetic laws in Maxwell's theory.
\end{remark}

\begin{definition}[Awareness]
Let $A \subseteq M$ denote the set of awareness states:
\begin{equation}
A \coloneqq \{ a \in M : a \text{ is directly available for global control and report} \}
\end{equation}
Awareness is the functional correlate of consciousness.
\end{definition}

\begin{principle}[Coherence Principle]\label{prin:coherence}
Consciousness and awareness systematically co-occur:
\begin{equation}
\forall x: \big( E(x) \neq \varnothing \big) \Leftrightarrow \big( A(x) \neq \varnothing \big)
\end{equation}
\end{principle}

\begin{principle}[Structural Coherence]\label{prin:structural_coherence}
The structure of phenomenal experience is isomorphic to the structure of awareness:
\begin{equation}
\boxed{S(E) \cong S(A)}
\end{equation}
where $S: \text{Domain} \to \text{Relational Structure}$ extracts the abstract relational/geometric structure.
\end{principle}

\begin{remark}
The isomorphism $S(E) \cong S(A)$ means:
\begin{enumerate}[nosep]
    \item The geometry of visual experience (e.g., 3D spatial relations) mirrors the structure of visual awareness representations.
    \item The structure of color space in experience (opponent relations, similarity metrics) mirrors the structure of color information in awareness.
    \item Temporal structure in experience mirrors functional temporal integration.
\end{enumerate}
\end{remark}

%------------------------------------------------------------------------------
\subsection{Invariance and Information Operators (Chapters 7 \& 8)}
%------------------------------------------------------------------------------

\begin{definition}[Fine-Grained Functional Organization]
For any system $x$, let $\Forg(x)$ denote its fine-grained functional organization:
\begin{equation}
\Forg(x) \coloneqq \langle Q_x, T_x, I_x, O_x \rangle
\end{equation}
where:
\begin{itemize}[nosep]
    \item $Q_x$ is the set of internal states
    \item $T_x: Q_x \times I_x \to Q_x$ is the state transition function
    \item $I_x$ is the input alphabet
    \item $O_x: Q_x \to O_x$ is the output function
\end{itemize}
Two systems are functionally isomorphic if their organizations are isomorphic as automata.
\end{definition}

\begin{principle}[Organizational Invariance]\label{prin:org_invariance}
Functional organization fully determines phenomenal experience:
\begin{equation}
\boxed{\forall x, y: \Forg(x) = \Forg(y) \Rightarrow E(x) \equiv E(y)}
\end{equation}
where $E(x) \equiv E(y)$ denotes qualitative identity of experience.
\end{principle}

\begin{theorem}[Fading Qualia Argument]
Absent qualia (functional isomorphs lacking experience) are naturally impossible.
\end{theorem}

\begin{proof}[Proof Sketch]
Consider gradual neural replacement by functional isomorphs.
\begin{enumerate}[nosep]
    \item If qualia faded gradually while function remained constant, the system would make erroneous judgments about its own experience (reporting ``bright red'' while experiencing ``faded pink'').
    \item Such systematic errors about direct phenomenal states are implausible.
    \item If qualia disappeared suddenly at some replacement threshold, this would be arbitrary.
    \item Therefore, qualia remain constant under functional isomorphism. \qed
\end{enumerate}
\end{proof}

\begin{definition}[Information Space]
Let $I$ be an information space:
\begin{equation}
I \coloneqq \{ s : s \text{ is a state selected from a space of alternatives} \}
\end{equation}
following Shannon's conception of information as ``difference that makes a difference.''
\end{definition}

\begin{principle}[Double-Aspect Principle of Information]\label{prin:double_aspect}
Information has two fundamental aspects---physical and phenomenal:
\begin{equation}
\boxed{\forall s \in I: \exists \PhiP(s) \in P \land \exists \PhiE(s) \in E}
\end{equation}
where:
\begin{itemize}[nosep]
    \item $\PhiP: I \to P$ is the physical realization function
    \item $\PhiE: I \to E$ is the phenomenal realization function
\end{itemize}
\end{principle}

\begin{remark}
The Double-Aspect Principle suggests that information is ontologically fundamental, with physical and phenomenal properties as dual manifestations:
\begin{equation}
\text{Physics} = \text{Information from the outside}; \quad \text{Experience} = \text{Information from the inside}
\end{equation}
\end{remark}

\begin{corollary}[Structure Preservation]
The realization functions preserve information structure:
\begin{equation}
S(I) \cong S(\PhiP(I)) \cong S(\PhiE(I))
\end{equation}
\end{corollary}

%------------------------------------------------------------------------------
\subsection{Quantum Integration (Chapter 10)}
%------------------------------------------------------------------------------

We extend the framework to quantum mechanics using the Everett interpretation.

\begin{definition}[Universal Quantum State]
The universal quantum state is a superposition of physical implementations:
\begin{equation}
\boxed{\ket{\Psiuniv} = \sum_{i} c_i \ket{P_i}}
\end{equation}
where $\{ \ket{P_i} \}$ is a basis of physical implementation states (e.g., pointer states or decoherence branches) and $c_i \in \mathbb{C}$ with $\sum_i |c_i|^2 = 1$.
\end{definition}

\begin{axiom}[Universal Schr\"odinger Evolution]
The only dynamical law is the Schr\"odinger equation:
\begin{equation}
i\hbar \frac{\partial}{\partial t} \ket{\Psiuniv} = \hat{H} \ket{\Psiuniv}
\end{equation}
There is no objective collapse postulate.
\end{axiom}

\begin{theorem}[Superposition of Minds]\label{thm:superposition_minds}
A superposed brain state implementing multiple distinct functional organizations yields multiple distinct phenomenal subjects:
\begin{equation}
\boxed{E\big( \ket{\Psiuniv} \big) \to \{ E_i \}_{i \in \mathcal{I}}}
\end{equation}
where $\{ E_i \}$ is a set of distinct phenomenal subjects, one for each decoherent branch.
\end{theorem}

\begin{proof}
\textbf{Step 1:} Let the universal state be $\ket{\Psiuniv} = \sum_i c_i \ket{P_i}$ where each $\ket{P_i}$ represents a distinct physical configuration.

\textbf{Step 2:} Each branch $\ket{P_i}$ instantiates a functional organization $\Forg(P_i)$.

\textbf{Step 3:} By the Principle of Organizational Invariance (Principle~\ref{prin:org_invariance}), each distinct functional organization determines a phenomenal state via $\Psi$:
\begin{equation}
E_i = \Psi(P_i) \quad \text{for each } i
\end{equation}

\textbf{Step 4:} If $\Forg(P_i) \neq \Forg(P_j)$, then $E_i$ and $E_j$ are distinct phenomenal subjects.

\textbf{Step 5:} The superposition does not yield a ``superposed mind'' but rather a multiplicity of definite minds:
\begin{equation}
E\big( \sum_i c_i \ket{P_i} \big) = \{ E_i : i \in \mathcal{I} \}
\end{equation}
where $\mathcal{I}$ indexes branches with distinct functional organizations. \qed
\end{proof}

\begin{remark}
This explains the phenomenology of definite experience despite physical indefiniteness: each observer experiences a definite outcome because each phenomenal subject $E_i$ is associated with a definite branch $\ket{P_i}$.
\end{remark}

\begin{corollary}[Resolution of the Measurement Problem]
The apparent ``collapse'' is explained without modifying quantum dynamics:
\begin{equation}
\text{Physical: } \ket{\Psiuniv} = \sum_i c_i \ket{P_i} \quad \longrightarrow \quad \text{Phenomenal: } \{ E_i \} \text{ (each definite)}
\end{equation}
The discreteness of experience emerges from the discreteness of functional organization, not from physical collapse.
\end{corollary}

%==============================================================================
\section{Summary: The Chalmers Ansatz}
%==============================================================================

The complete formal framework is:

\begin{tcolorbox}[colback=yellow!5,colframe=orange!70,title=\textbf{The Chalmers Ansatz}]
\textbf{Universe Tuple:}
\begin{equation}
\mathcal{U} = \langle P, M, E, L \rangle \quad \text{with } L = L_P \cup L_\Psi
\end{equation}

\textbf{Supervenience Relations:}
\begin{align}
P &\logent M && \text{(Psychological: logical supervenience)} \\
P &\notlogent E && \text{(Phenomenal: logical supervenience \textbf{fails})} \\
P &\natent E && \text{(Phenomenal: natural supervenience holds)}
\end{align}

\textbf{Psychophysical Laws:}
\begin{align}
\Psi &: P \to E && \text{(Bridging function)} \\
S(E) &\cong S(A) && \text{(Structural coherence)}
\end{align}

\textbf{Invariance Principles:}
\begin{align}
\Forg(x) = \Forg(y) &\Rightarrow E(x) \equiv E(y) && \text{(Organizational invariance)} \\
\forall s \in I: \exists \PhiP(s) &\land \exists \PhiE(s) && \text{(Double-aspect information)}
\end{align}

\textbf{Quantum Extension:}
\begin{equation}
\ket{\Psiuniv} = \sum_i c_i \ket{P_i} \quad \longrightarrow \quad E(\ket{\Psiuniv}) = \{ E_i \}
\end{equation}
\end{tcolorbox}

\vfill

\begin{center}
\rule{0.5\textwidth}{0.4pt}\\[0.5em]
\small Generated using K-Dense Web (\href{https://k-dense.ai}{k-dense.ai})
\end{center}

\end{document}
